\documentclass[aip,reprint,floatfix,amsmath,amssymb]{revtex4-2}
\usepackage{graphicx}
\usepackage{booktabs}
\usepackage{array}
\usepackage{multirow}
\usepackage{makecell}
\usepackage{siunitx}
\usepackage{xcolor}
\usepackage{hyperref}
\usepackage{url}
\hypersetup{colorlinks=true,linkcolor=blue!45!black,citecolor=blue!45!black,urlcolor=blue!45!black,pdftitle={Reference Information and Target-Power Regimes in Multi-Window Photovoltaic Forecasting},pdfauthor={Jiangkun Zhu et al.}}
\sisetup{detect-all,round-mode=places,round-precision=3}
\setcounter{dbltopnumber}{4}\renewcommand{\dbltopfraction}{0.95}\renewcommand{\textfraction}{0.05}

\begin{document}
\preprint{WORKING MANUSCRIPT --- AUTHOR DECLARATIONS PENDING CONFIRMATION}
\title{Reference Information and Target-Power Regimes in Multi-Window Photovoltaic Forecasting}

\author{Jiangkun Zhu}
\thanks{ORCID: \href{https://orcid.org/0009-0009-5335-2345}{0009-0009-5335-2345}}
\affiliation{College of Advanced Manufacturing, Fuzhou University, 1 Shuicheng Road, Quanzhou 362251, China}
\author{Mengling Yang}
\affiliation{College of Advanced Manufacturing, Fuzhou University, 1 Shuicheng Road, Quanzhou 362251, China}
\author{Zhicong Chen}
\email{zhicong.chen@fzu.edu.cn}
\affiliation{College of Advanced Manufacturing, Fuzhou University, 1 Shuicheng Road, Quanzhou 362251, China}
\affiliation{College of Physics and Information Engineering, Fuzhou University, 2 Xueyuan Road, Fuzhou 350116, China}
\author{Lijun Wu}
\email{lijun.wu@fzu.edu.cn}
\affiliation{College of Advanced Manufacturing, Fuzhou University, 1 Shuicheng Road, Quanzhou 362251, China}
\affiliation{College of Physics and Information Engineering, Fuzhou University, 2 Xueyuan Road, Fuzhou 350116, China}



\begin{abstract}
Photovoltaic forecasting gains depend on the reference information and the target-power range being evaluated. We examine these dependencies in five power systems at three geographic sites: three co-located Alice Springs arrays and independently fitted Yulara and NIST Ground replications. Four compact neural implementations predict twelve-hour trajectories from six-hour histories in 60 historical training runs. Cumulative errors use the first 1, 4, 8, or 12 hours, with matched targets for each reference. Daily Persistence outperforms the post hoc neural envelope in 22 of 24 Alice comparisons. The prespecified external Inverted-variate model improves on Last-value in 16 of 16 comparisons and on Daily in 15, while neural rankings differ at NIST. Its NIST twelve-hour full-target error nevertheless exceeds Daily by 6.47\%: a 420.64 kW$^2$ reduction in the power-active contribution to mean squared error is offset by a 655.38 kW$^2$ increase in the low-power contribution. Post hoc linear references further show that adding available previous-day trajectories lowers twelve-hour Ridge error in all five systems, without consistently outperforming Daily or the neural model. Temporal-block intervals qualify the small long-window gains. The results identify reference history, target-power composition, and evaluation support as essential parts of a reported forecasting advantage.
\end{abstract}
\keywords{photovoltaic forecasting, reference information, cumulative forecast windows, error decomposition}
\maketitle

\section{Introduction}
Photovoltaic (PV) forecasting supports storage scheduling and power-system operation by estimating how generation will change after a forecast is issued. The practical question is what a learned forecast adds to information already available at that time. Repeating the latest power and retrieving yesterday's corresponding power supply different references. A gain over the former need not imply a gain over the latter. Moreover, aggregate error combines periods of substantial production with low-output targets, whose frequencies and errors can change the comparison.\cite{Antonanzas2016,Mellit2021Review,Hyndman2006}

We study these dependencies through a co-located comparison at Alice Springs and two independently fitted external-site replications. Five target systems cover three geographic sites, with separate models and preprocessing for each system. Four compact neural implementations use six hours of observations to predict a twelve-hour trajectory. The original neural experiment is supplemented by prediction-only analyses and two deterministic linear references: Ridge using the same recent history, and Ridge augmented with an available previous-day trajectory. These additions are post hoc analyses of performance splits already observed during the project.

The principal finding concerns the interpretation of forecasting gains. Daily Persistence changes the conclusion drawn from Last-value Persistence at Alice Springs; the external primary model usually improves on both, but its NIST twelve-hour full-target comparison reverses despite lower power-active error. An additive squared-error decomposition quantifies that reversal. Common-origin and temporal-block analyses then separate differences in target support and training randomness from variation over the observed dates. The linear references test whether a simple learned mapping and added daily history change the comparison, without extending the neural architecture search.

\section{Related work}
PV forecasting research uses historical measurements, physical models, numerical weather prediction, sky imagery, and their combinations. These information sets support different forecasting tasks.\cite{Antonanzas2016,Mellit2021Review,DiLeo2025Review} Weather-based benchmarks compare classical and deep models under future weather inputs; geographically distributed studies also investigate transfer.\cite{Markovics2026,SolNet2025} Their numerical errors are not interchangeable with those of a history-only, separately fitted experiment. Studies using Alice Springs and Yulara establish the value of these archives, while differing in periods, available fields, and reference definitions.\cite{Mahmud2021,Huang2024MANODE,Mansour2026}

Recurrent, variate-token, patch-token, and convolutional models provide alternative ways to summarize a measured history.\cite{Hewamalage2021RNN,Liu2024iTransformer,Nie2023PatchTST,Luo2024ModernTCN} Here these mechanisms are represented by compact project implementations with a fixed training budget. Their comparison addresses the implemented forecasters, rather than optimally tuned official architectures. Forecast-evaluation guidance motivates chronological splits, Train-only transformations, and explicit sample support.\cite{Tashman2000,Bergmeir2012,Hewamalage2023,Kapoor2023} We use these controls as the basis for interpreting reference-dependent error, rather than as an algorithmic contribution.

\section{Study design and information sets}
\subsection{Five targets at three sites}
Figure~\ref{fig:protocol} shows the information available at forecast time. Alice Springs contributes three co-located arrays and 36 historical neural runs; Yulara and NIST Ground each contribute one target and 12 runs. Seeds are 42, 43, and 44. No observations are pooled across systems. Table~\ref{tab:data} records the two input regimes and chronological splits.

\begin{figure*}[!tp]
\centering\includegraphics[width=0.97\textwidth]{review_figures/fig1_design.pdf}
\caption{Independent fitting and forecast-time information. Every step is five minutes; cumulative windows use the first 12, 48, 96 or 144 future steps.}
\label{fig:protocol}
\end{figure*}
\begin{table*}[!tp]
\caption{Three geographic sites and two input regimes. Capacities are descriptive DC ratings; error normalization uses each target's Train range. All forecast grids are five-minute grids. Dates in the last two rows refer to 2017, and those in the first row to 2018.}
\label{tab:data}\centering\small\setlength{\tabcolsep}{3pt}
\begin{tabular}{lllll}
\toprule Site / systems & \makecell[l]{DC rating /\\channels} & \makecell[l]{Train; Validation;\\Test} & \makecell[l]{Source resolution /\\time basis} & Interpretation limit\\\midrule
\makecell[l]{Alice Springs /\\three co-located\\arrays} & \makecell[l]{Sanyo 6.30, Hanwha\\5.83, Qcells 5.90 kW\\17 channels} & \makecell[l]{Apr 1--Jul 15;\\Jul 16--Aug 7;\\Aug 8--31} & \makecell[l]{5 min / frozen\\ACST coordinate} & \makecell[l]{Shared weather;\\technology is not\\an isolated treatment}\\
\makecell[l]{Yulara / one\\combined facility} & 106.6 kW / 7 & \makecell[l]{Jan--Aug; Sep;\\Oct--Dec} & \makecell[l]{5 min / provider local;\\+5 min availability} & \makecell[l]{GHI metadata and UTC\\offset partly\\unconfirmed}\\
\makecell[l]{NIST Ground /\\one array} & 271 kW / 7 & \makecell[l]{Jan--Aug; Sep;\\Oct--Dec} & \makecell[l]{1 min to 5 min /\\fixed EST, UTC$-05{:}00$,\\no DST} & \makecell[l]{Complete 2017\\operation log\\unavailable}\\
\bottomrule\end{tabular}
\end{table*}
The Alice targets are Sanyo hybrid-silicon, Hanwha polycrystalline-silicon, and Qcells monocrystalline-silicon arrays.\cite{DKASC2026,Wang2019DKASC} Eight numerical variables comprise array power and Performance Ratio, ambient temperature, relative humidity, and global/diffuse horizontal/tilted radiation. Eight original missingness masks and one Isolation Forest flag bring the input width to 17. Co-location provides shared weather context; differences between arrays also include their electrical systems, orientation, and sensing. Technology is therefore not an isolated experimental treatment.

The Yulara target is a 106.6 kW combined facility, not multiple independent sites.\cite{YularaData} Both external systems use seven channels in a fixed order: historical AC power, ambient temperature, global horizontal irradiance (GHI), the three corresponding missingness masks, and an Isolation Forest flag. This common available information set is fitted independently at each site. It does not transfer the original 17-channel weights. Yulara's temperature is used as the provider's Celsius field and GHI in its native numerical units, with no additional conversion or cross-site absolute irradiance inference.

NIST power comes from the AC real-power meter; its inverter-power field, including the $-999$ code, is excluded. Temperature and GHI use the official averaged fields, with the supplied W/m$^2$ field used directly.\cite{Boyd2017NIST,NISTDictionary,NISTData} The fixed EST/LST coordinate retains UTC$-05{:}00$ throughout, without daylight-saving changes. At each five-minute endpoint $T$, only observations in $[T-5\mathrm{\ min},T)$ contribute. A variable requires five distinct finite minute values; otherwise its aggregate is missing. Yulara retains the provider-local coordinate with a conservative five-minute availability delay. Two off-grid records are excluded without rounding. These availability conventions prevent a completed measurement interval from entering an earlier input.

Finite negative power and GHI, including unusual operating conditions, remain in the external records. The original Alice rule retains finite nonnegative labels. Missing labels are never filled. Neither low output nor a power--GHI discrepancy is used to remove Test records. The Supplement gives field definitions, time audits, and the unavailable NIST operation log and partial Yulara metadata.

\subsection{Preprocessing and forecast definition}
Original numerical missingness is recorded before input filling. A five-neighbor KNN imputer is fitted on Train. Isolation Forest then uses the imputed Train values (100 trees, contamination 0.01, random state 42); its flag and the original masks are appended. A feature MinMaxScaler and a separate target scaler are fitted on Train only, the latter using valid unimputed power.\cite{Troyanskaya2001,Liu2008Isolation} Validation and Test use the fitted transformations unchanged. Missing timeline positions remain in place.

At origin $t$, the 72 input positions end at $t$, and their values must be available by $t$. The first target is $t+5$ min and the final target $t+12$ h. Inputs and labels remain within one split. H12, H48, H96, and H144 denote \emph{cumulative forecast windows}: the first 12, 48, 96, and 144 targets, corresponding to 1, 4, 8, and 12 h. They are not errors at four isolated leads. Original Test evaluation retains origins with valid labels throughout the requested prefix, so different windows can have different origins. A separate common-origin analysis uses complete H144 origins for every prefix; fixed-lead curves use these same origins and score one lead at a time.

\section{Forecasters and matched evaluation}
\subsection{Compact neural implementations}
The recurrent forecaster combines multi-dilation convolutions, an attention-encoder branch, gated fusion, and an attention-weighted bidirectional GRU summary. Its trajectory decoder obeys
\begin{equation}
s_k=\operatorname{GRUCell}(c,s_{k-1}),\qquad \hat y_k=\operatorname{readout}(s_k),\quad s_0=c.
\end{equation}
The fixed context $c$ is supplied at every step. Previous predicted power is not fed back, so this is recurrent hidden-state trajectory decoding, rather than power-autoregressive forecasting. The other three implementations use direct multi-output heads. Inverted-variate embeds each channel's complete history as a token; Joint-patch embeds overlapping multivariate patches; Depthwise TCN uses an input projection and depthwise/pointwise temporal blocks. Layer specifications and parameter counts for both input widths are in the Supplement. All reported rankings refer to these compact implementations.

Historical neural training used batch size 256, AdamW with learning rate $10^{-3}$ and weight decay $10^{-5}$, gradient clipping 1.0, a maximum of 25 epochs, and early-stopping patience five with minimum improvement $10^{-8}$. Training and checkpoint selection require complete H144 targets. The selected checkpoint minimizes Validation total squared error divided by total valid-target count, not an unweighted average of batch losses. This common budget is a controlled comparison, not evidence that every architecture is fully optimized. In particular, all three NIST TCN runs have their best Validation checkpoint at epoch 25; the historical learning curves reach the imposed boundary.

\subsection{References and target-power scopes}
Last-value repeats the true valid power at the origin. Daily Persistence retrieves power at each target timestamp minus exactly 24 h through a timestamp join. At the longest lead, this observation is still twelve hours before the origin. Daily thus uses a different history from the six-hour neural input. On Daily-matched comparisons, every method uses the same origins, labels, and elementwise intersection of valid target and Daily-lag points. This ensures identical scored targets, not identical information.

The full scope includes all eligible targets. The \emph{power-active} scope uses true future power greater than 1\% of that system's maximum Train power. This is a descriptive, post hoc evaluation mask, not a forecast-time detector. Its complement is \emph{low-power}; it may include low production during daytime. The mask does not identify night, snow, maintenance, or outages. Frozen machine-readable records retain the historical key \texttt{daylight}; its numerical definition and results are unchanged.

\subsection{Metrics and analysis hierarchy}
RMSE, MAE, bias, and $R^2$ use the stated common targets. Train-range nRMSE divides RMSE by that system's Train power range. Reported nRMSE and skill use percentages, with skill $100(1-\mathrm{RMSE}_{m}/\mathrm{RMSE}_{r})$; a zero reference error leaves skill undefined. Negative $R^2$ is retained, and zero target variance leaves it undefined. Neural summaries average per-seed metrics, with sample SD ($\mathrm{ddof}=1$), rather than evaluating a mean-prediction ensemble. Baselines are deterministic. Seed SD describes training randomness, not temporal sampling uncertainty.

Inverted-variate was fixed as the external primary model from the frozen Alice comparison before external Test predictions were generated. It is a descriptive focal model at Alice. Four-model rankings are secondary. The post hoc minimum mean neural RMSE is a descriptive envelope favorable to the neural group, not a prespecified or deployable forecast. Win counts across windows and scopes are dependent descriptive summaries. We do not pool cross-site kW errors or add the Alice 24 and external 16 combinations into an overall win total.

\subsection{Post hoc explanatory analyses and linear references}
The additional analysis plan was recorded before calculating its new scores, after the original Test results were known. The power decomposition uses disjoint active and low-power targets:
\begin{equation}
\mathrm{MSE}_{\rm full}=\frac{N_a}{N}\mathrm{MSE}_a+\frac{N_l}{N}\mathrm{MSE}_l.
\end{equation}
Contributions are computed per seed and then averaged. Common-origin sensitivity and fixed-lead errors distinguish window length from changing sample composition. Monthly external summaries assign complete trajectories to the forecast-origin month. Calendar cases use the first legal origin on or after the 15th at 10:00 in each Test month; additional best/worst cases are explicitly post hoc diagnostics, always using seed 42.

For temporal uncertainty, paired nonoverlapping 48 h origin blocks are sampled with replacement 2000 times, with 24/72 h sensitivity checks and fixed random seed 20260908. The same block indices are used for methods and neural seeds. Each resample sums squared errors and counts before computing per-seed RMSE and skill, then averages the seed skills. Percentile 95\% intervals are conditional on observed dates and fitted models. Blocks preserve local dependence but cannot remove all dependence between overlapping trajectories across block boundaries; the short Alice Test period provides especially limited temporal support.

Ridge A flattens the corresponding 72-step, 17- or seven-channel processed neural history. Ridge B adds 144 exact previous-day target-time power values and their missingness masks. Missing lag inputs use Train-column medians. Standardization is Train-fitted, labels are not imputed, and both variants use complete H144 Train/Validation windows. One common regularization parameter from $\{0.1,1,10,100,1000\}$ minimizes full H144 Validation squared error, with larger regularization breaking an exact tie. The solver is deterministic and uses no seed repetitions. Test uses the original horizon-specific origins and common Daily-valid masks. These references test recent-history learning and additional daily information; they do not predict the effect of giving neural models the same added history.

\section{Results}
\subsection{Alice Springs: the reference changes the comparison}
The frozen primary winners remain Inverted-variate 12, TCN 9, Joint-patch 2, recurrent 1, and Last-value 0 across 24 array--window--scope combinations. Table~\ref{tab:primarysummary} retains the short and full-window primary summary. The neural envelope improves on Last-value throughout, but Daily has lower matched error than that favorable envelope in 22 of 24 combinations. The two neural exceptions are Hanwha's one-hour full and power-active comparisons. Figure~\ref{fig:alice} places both references on the same Daily-valid targets, so their contrast cannot be explained by using different target subsets for the two baselines.

% AUTO-GENERATED by build_figures.py from corrected_metrics.csv. DO NOT EDIT.
\begin{table*}[!tp]
\caption{Primary Test Train-range nRMSE and skill (\%) on horizon-specific valid origins. The envelope is the post hoc minimum three-seed mean among the four compact implementations and is a descriptive upper bound, not a prespecified model; the complete model-by-seed table is in the Supplementary Material.}
\label{tab:primarysummary}\centering\small
\begin{tabular}{lllrrrr}
\toprule Array & Scope & Window & Envelope member & Envelope nRMSE & Last-value nRMSE & RMSE skill\\\midrule
Sanyo & Full & 1 h & Discrete recurrent & 3.42 & 7.66 & 55.30\\
Sanyo & Full & 12 h & Inverted-variate & 12.80 & 45.06 & 71.60\\
Sanyo & Power-active & 1 h & Depthwise TCN & 4.80 & 11.01 & 56.40\\
Sanyo & Power-active & 12 h & Inverted-variate & 17.45 & 60.56 & 71.20\\
Hanwha & Full & 1 h & Depthwise TCN & 3.01 & 7.45 & 59.50\\
Hanwha & Full & 12 h & Depthwise TCN & 12.75 & 43.94 & 71.00\\
Hanwha & Power-active & 1 h & Depthwise TCN & 3.96 & 10.78 & 63.30\\
Hanwha & Power-active & 12 h & Depthwise TCN & 17.82 & 56.14 & 68.30\\
Qcells & Full & 1 h & Inverted-variate & 5.42 & 7.80 & 30.50\\
Qcells & Full & 12 h & Depthwise TCN & 13.07 & 40.05 & 67.40\\
Qcells & Power-active & 1 h & Inverted-variate & 7.56 & 11.29 & 33.00\\
Qcells & Power-active & 12 h & Depthwise TCN & 19.09 & 60.68 & 68.50\\
\bottomrule\end{tabular}\end{table*}


\begin{figure*}[!tp]
\centering\includegraphics[width=0.97\textwidth]{review_figures/fig2_alice_references.pdf}
\caption{Post hoc minimum three-seed mean neural RMSE relative to both references on the same Daily-valid targets. F: full; A: power-active. The envelope is descriptive, not deployable.}
\label{fig:alice}
\end{figure*}

Horizon-specific eligibility matters particularly for Qcells: its one-hour primary evaluation has 6463 origins and 77556 full targets, including 36504 power-active targets, whereas complete H144 support has 2996 origins. The original short-window result is preserved. Common-origin and fixed-lead curves in the Supplement describe the retained complete trajectories and do not recover the missing periods excluded from that set. Consequently, a change along the original cumulative-window curve cannot be attributed solely to a longer forecast window. For Qcells at one hour, common-H144 support contains 35952 full targets but only 42 power-active points, compared with 36504 active points in the original analysis. On that selected full set, Last-value RMSE is 0.00398 kW versus 0.12494 kW for Inverted-variate, and Last-value leads all four neural models. This sensitivity exposes a major composition change rather than invalidating the original horizon-specific comparison.

\subsection{External replications: recurrent gains, conditional rankings}
The external prespecified model has lower mean RMSE than Last-value in all 16 comparisons and than Daily in 15. Table~\ref{tab:external} and Fig.~\ref{fig:external} retain every comparison and the three-seed dispersion. At twelve hours/full, Last-value skill is 67.10\% at Yulara and 43.24\% at NIST; Daily-matched skill is 1.62\% and $-6.47$\%, respectively. The reference therefore changes both the size and, at NIST, the sign of the observed advantage.

\begin{table*}[!tp]
\centering\small
\caption{External prespecified Inverted-variate results: three-seed mean $\pm$ sample SD. nRMSE and skill are percentages; Daily skill uses its own common target intersection.}
\label{tab:external}
\begin{tabular}{lllrrr}
\toprule Site & Window & Scope & Train-range nRMSE & Last-value skill & Daily-matched skill\\\midrule
Yulara & 1 h & full & $9.36\pm0.03$ & $17.34\pm0.26$ & $36.28\pm0.20$\\
Yulara & 1 h & power-active & $12.72\pm0.06$ & $18.73\pm0.36$ & $37.28\pm0.28$\\
Yulara & 4 h & full & $12.16\pm0.13$ & $49.24\pm0.53$ & $17.25\pm0.86$\\
Yulara & 4 h & power-active & $16.06\pm0.22$ & $48.58\pm0.70$ & $20.83\pm1.08$\\
Yulara & 8 h & full & $14.28\pm0.05$ & $61.32\pm0.14$ & $2.71\pm0.34$\\
Yulara & 8 h & power-active & $19.03\pm0.14$ & $56.10\pm0.31$ & $6.10\pm0.67$\\
Yulara & 12 h & full & $14.37\pm0.14$ & $67.10\pm0.32$ & $1.62\pm0.97$\\
Yulara & 12 h & power-active & $19.28\pm0.20$ & $59.46\pm0.42$ & $4.50\pm0.99$\\
NIST Ground & 1 h & full & $6.91\pm0.04$ & $17.93\pm0.52$ & $57.73\pm0.28$\\
NIST Ground & 1 h & power-active & $11.01\pm0.04$ & $20.24\pm0.26$ & $59.01\pm0.13$\\
NIST Ground & 4 h & full & $11.05\pm0.10$ & $41.16\pm0.51$ & $32.09\pm0.58$\\
NIST Ground & 4 h & power-active & $17.22\pm0.39$ & $37.98\pm1.39$ & $36.07\pm1.41$\\
NIST Ground & 8 h & full & $15.59\pm0.29$ & $42.18\pm1.06$ & $3.88\pm1.80$\\
NIST Ground & 8 h & power-active & $21.58\pm0.30$ & $39.25\pm0.85$ & $20.22\pm1.12$\\
NIST Ground & 12 h & full & $17.29\pm0.27$ & $43.24\pm0.87$ & $-6.47\pm1.65$\\
NIST Ground & 12 h & power-active & $23.49\pm0.40$ & $37.80\pm1.06$ & $13.20\pm1.49$\\
\bottomrule\end{tabular}
\end{table*}

\begin{figure*}[!tp]
\centering\includegraphics[width=0.97\textwidth]{review_figures/fig3_external_skills.pdf}
\caption{Prespecified external Inverted-variate skill. Dots are individual seeds; bars show sample SD of three seed metrics, not temporal confidence intervals. Each reference uses its own matched support.}
\label{fig:external}
\end{figure*}

Inverted-variate leads all eight Yulara neural comparisons. At NIST its mean neural rank is 2.25: recurrent decoding leads the full scope, while TCN leads the longer power-active windows. Figure~\ref{fig:gaps} shows error gaps from the best neural mean instead of amplifying small differences through ranks alone. The post hoc neural envelope beats Daily in eight Yulara and seven NIST comparisons, preserving NIST twelve-hour/full as the Daily exception. These rankings describe the fixed budget and do not establish a site-invariant ordering of architecture families.

\begin{figure*}[!tp]
\centering\includegraphics[width=0.97\textwidth]{review_figures/fig4_neural_gaps.pdf}
\caption{Relative RMSE gaps from the best mean neural implementation in each cumulative window and power scope. Zero identifies the lowest error; small gaps are shown as small gaps rather than enlarged rank differences.}
\label{fig:gaps}
\end{figure*}

\subsection{Explaining the NIST full-target reversal}
NIST twelve-hour Daily-matched support contains 3128676 target points, of which 1098178 are power-active and 2030498 (64.90\%) are low-power. Inverted-variate's mean per-seed RMSE is 44.61 kW overall, 60.54 kW active, and 32.91 kW low-power. Daily's corresponding errors are 41.90, 69.75, and 8.60 kW. The active subset favors the neural forecast, while the larger low-power subset carries a substantial extra error.

Table~\ref{tab:decomp} and Fig.~\ref{fig:decomp} quantify this cancellation on the squared-error scale. The neural model reduces the active contribution to full MSE by 420.64 kW$^2$ but increases the low-power contribution by 655.38 kW$^2$, leaving an overall increase of 234.74 kW$^2$. The neural low-power bias is positive (12.73 kW), compared with 1.11 kW for Daily. This directly explains the aggregate reversal through scored target subsets; it does not identify a weather event or a network mechanism. The decomposition is also supplied for every system, model, and window.
\begin{table}[!tp]\caption{NIST twelve-hour contributions to full MSE (kW$^2$). Contributions, not RMSEs, add across disjoint target subsets. Neural values average per-seed contributions.}\label{tab:decomp}\centering\small\begin{tabular}{lrrr}\toprule Method & Active & Low-power & Total\\\midrule
Inverted & 1286.91 & 703.40 & 1990.31\\
Daily & 1707.55 & 48.02 & 1755.56\\
\bottomrule\end{tabular}\end{table}

\begin{figure*}[!tp]
\centering\includegraphics[width=0.97\textwidth]{review_figures/fig5_decomposition_case.pdf}
\caption{Left: additive active/low-power contributions to full MSE on NIST twelve-hour Daily-matched targets; neural contributions are averaged over seeds. Right: one seed42 trajectory from the fixed calendar selection, not a rolling reconstruction.}
\label{fig:decomp}
\end{figure*}

\subsection{Temporal support and simple learned references}
Paired 48 h block intervals qualify the small long-window Daily gains. Yulara twelve-hour/full skill is 1.62\%, with a conditional interval of $[-9.47,9.85]$\%; NIST eight-hour/full is 3.88\% with $[-6.10,11.91]$\%, and twelve-hour/full is $-6.47$\% with $[-17.14,2.28]$\%. NIST's corresponding power-active eight- and twelve-hour intervals are $[11.75,27.07]$\% and $[4.60,19.94]$\%. All prechosen block-size results are retained in the Supplement. These intervals describe temporal variation over the observed Test support and fitted models, separately from seed SD.

Table~\ref{tab:ridge} shows the post hoc twelve-hour linear comparison. Adding daily-history inputs lowers Ridge error in all five systems. At NIST, Ridge B reaches 35.00 kW, compared with Ridge A's 48.96 kW, Daily's 41.90 kW, and Inverted-variate's 44.61 kW. At Yulara, both linear variants remain above Daily and the neural model; the large error of Ridge A is retained despite its Validation-selected regularization. Across Alice arrays, daily-augmented Ridge is much closer to Daily, exceeding it at Sanyo and Hanwha and slightly improving at Qcells. Thus simple learning alone does not explain all observed neural gains, while available daily information can materially change a learned reference.
\begin{table}[!tp]\caption{Twelve-hour Daily-matched full-target RMSE (kW). Neural values are means of seed metrics; Ridge references are deterministic and post hoc.}\label{tab:ridge}\centering\small\begin{tabular}{lrrrr}\toprule System & Inverted & Daily & Ridge A & Ridge B\\\midrule
Sanyo & 0.775 & 0.118 & 0.914 & 0.212\\
Hanwha & 0.805 & 0.203 & 0.844 & 0.248\\
Qcells & 0.844 & 0.244 & 0.839 & 0.237\\
Yulara & 15.805 & 16.066 & 119.585 & 35.950\\
NIST & 44.609 & 41.899 & 48.956 & 34.995\\
\bottomrule\end{tabular}\end{table}

\section{Interpretation and limits}
Last-value holds the present power through the entire future trajectory, whereas learned recent-history models can represent changing output. The repeated neural improvement establishes added predictive value over that local-continuity reference. Daily supplies a measured trajectory from another day and therefore asks whether a model adds to daily recurrence. The Alice--external contrast and the daily-augmented Ridge results show why these questions should be reported separately. Differences in daily stability, operating distributions, and input availability are possible explanations for site differences, not mechanisms established by this experiment.

The power decomposition provides a direct explanation for one scope-dependent result: NIST's active error saving is outweighed by low-power error. It also explains why the full-target winner need not be the active-target winner. This conclusion follows from disjoint squared-error contributions and does not require attributing the low-power records to night, snow, or downtime. Ranking differences across sites remain observational, with different years, input widths, target systems, and held-out months.

The three geographic sites broaden the evidence beyond a single facility, but the external October--December period is not an all-season validation. Alice's three arrays are co-located and use 2018 data with 17 channels, compared with seven external channels in 2017. Each model is fitted separately. The missing NIST operation log and partial Yulara metadata limit event and sensor interpretation. Test records underwent data-quality inspection during development; they remain held-out performance splits, while the added decompositions, intervals, and Ridge comparisons are explicitly post hoc. Fixed-budget compact models do not represent fully tuned official implementations. Seed SD and finite-period block intervals address different, incomplete sources of uncertainty. Common target masks remove within-comparison sample discrepancies, not information-set differences or selection effects. The Qcells one-hour common-origin reversal demonstrates why the retained target distribution must also be reported.

\section{Conclusion}
The benchmark identifies three conditions that change a PV forecasting advantage: the reference history, the cumulative window and its support, and the target-power composition. Daily outperforms the favorable neural envelope in 22 of 24 Alice comparisons; the external prespecified model improves on Last-value in 16 of 16 and Daily in 15 of 16. At NIST's twelve-hour full scope, a quantified low-power error contribution reverses the active-target advantage. Added previous-day inputs improve a simple linear reference in all five systems without producing a uniform winner. Reporting these reference and support conditions makes a forecasting gain interpretable beyond an architecture name.
\section*{Supplementary Material}
The Supplement adds common-origin and fixed-lead analyses, paired temporal-block intervals, target-power decompositions, deterministic Ridge references, monthly summaries and transparent trajectory cases alongside complete Alice evidence, external field and time audits, all 96 support groups, the 24-run matrix and checkpoint records, per-seed RMSE, complete mean/sample-SD errors, metadata limits, model configurations, and the descriptive envelope. Full per-seed metric CSVs accompany the code rather than being expanded into thousands of PDF rows.

\section*{Acknowledgments}
The authors acknowledge the DKA Solar Centre and NIST for providing the source photovoltaic and meteorological data.
% FINAL_AUTHOR_CONFIRMATION_REQUIRED: declaration drafts are retained from the
% authors' prior materials and must be approved for this manuscript before upload.
\section*{Funding}
\textbf{Author confirmation required.} Draft: This work is supported by the National Natural Science Foundation of China (Grant Nos. 62271151 and W2421092).

\section*{Conflict of Interest}
\textbf{Author confirmation required.} Draft: The authors declare that they have no known competing financial interests or personal relationships that could have appeared to influence the work reported in this paper.

\section*{Ethics Approval}
\textbf{Author confirmation required.} Draft: This study analyzes publicly accessible photovoltaic and meteorological records and involves no human participants or animals; ethics approval was not applicable.

\section*{Author Contributions}
\textbf{Author confirmation required.} Draft: Jiangkun Zhu: Conceptualization, Methodology, Software, Validation, Formal analysis, Investigation, Data curation, Writing -- Original Draft, Visualization. Mengling Yang: Validation, Writing -- Review \& Editing. Zhicong Chen: Conceptualization, Supervision, Project administration, Funding acquisition, Writing -- Review \& Editing. Lijun Wu: Supervision, Resources, Project administration, Funding acquisition, Writing -- Review \& Editing.

% FINAL_AUTHOR_CONFIRMATION_REQUIRED: exact AI tool name, version, provider,
% purpose, and wording must be confirmed against AIP policy before upload.
\section*{Declaration of AI-Assisted Tools}
\textbf{Author confirmation required.} The exact tool/version record and any journal-required disclosure wording remain pending author confirmation; no model designation is inferred.

\section*{Data Availability}
\textbf{Author confirmation required.} Draft: Source observations are available from the official DKA Solar Centre services for Alice Springs and Yulara and from the NIST PV data service.\cite{DKASC2026,YularaData,NISTData} Reproduction requires downloading data from those providers and complying with their terms. The code repository does not redistribute these raw observations. The final public code location and release details remain subject to author confirmation.

\section*{Code Availability}
% FINAL_AUTHOR_CONFIRMATION_REQUIRED: replace with the released repository URL.
\textbf{Author confirmation required.} Draft: A dedicated public release of evaluation code, aggregate metrics, figure generation, and ordinary protocol tests is planned after manuscript approval and license selection. Raw data, checkpoints, and prediction arrays will not be redistributed. The final repository URL must be confirmed before submission.

\bibliography{references}
\end{document}
